\section{Results and Discussion}
\label{sec:results}

The methodology was applied end-to-end to the conceptual design of a
small electric tail-sitter VTOL UAV inspired by early V-BAT-class
demonstrators; the complete design repository --- configuration,
physics library, notebooks, decision records, and full commit history
--- is openly available \citep{tuncer2026vbat}. The configuration
combines a single ducted fan in the aft fuselage, four jet vanes in
the fan efflux providing pitch/yaw/roll authority in hover and
transition, wing ailerons as the cruise-phase roll effector, and
tail-down vertical takeoff and landing. The mission requirement is a
0.5\,kg payload carried through a deliberately short-hover profile ---
120\,s of vertical flight, 40\,s of transitions (billed at hover
power), and 900\,s of cruise at 20\,m/s ($\approx$18\,km range) with a
20\% energy reserve. This section reports (i) the converged design
outcome, (ii) the design-point evolution that produced it, (iii)
process metrics extracted from the project's git history quantifying
the AI's role, and (iv) discussion against the research question.

\subsection{Converged design point}
\label{sec:design-results}

Closure is performed with transparent first-order physics: hover power
from momentum (actuator-disk) theory with an overall figure-of-merit
efficiency chain, cruise power from a conservative lift-to-drag model,
and battery mass from mission energy through pack-level specific
energy, with the pack's Joule losses folded in self-consistently (the
battery efficiency is the mission-averaged $I^2R$ efficiency of the
nominal 90\,m$\Omega$ pack, iterated to the closure fixed point).
Table~\ref{tab:designpoint} summarizes the converged vehicle;
Fig.~\ref{fig:vehicle} shows the parametric solid model regenerated by
the pipeline on every run.

\begin{table}[htbp]
\centering
\caption{Converged conceptual design point (v0.5.2 snapshot).}
\label{tab:designpoint}
\footnotesize
\setlength{\tabcolsep}{4pt}
\begin{tabular}{@{}llll@{}}
\toprule
\textbf{Quantity} & \textbf{Value} & \textbf{Quantity} & \textbf{Value} \\
\midrule
MTOW (closure) & 2.518\,kg & Wing loading & 131.2\,N/m$^2$ \\
All-up mass (as-selected) & 2.336\,kg & Wing area / span & 0.1883\,m$^2$ / 1.063\,m \\
Payload & 0.5\,kg & Aspect ratio & 6.0 \\
Hover electrical power & 743\,W & Airfoil & NACA 4412 ($C_{L,\max,3D}=1.489$) \\
Peak battery discharge & 8.1\,C & Cruise $C_L$ / $L/D$ & 0.535 / 13.35 \\
Rotor & 203\,mm, 3-blade $8\times6$ & Cruise speed & 20\,m/s ($\approx$1.02\,$V_{\mathrm{MD}}$) \\
Blade section / solidity & Clark~Y / $\sigma=0.077$ & Mission & 120\,s hover + 40\,s trans.\ + 900\,s cruise \\
1/rev forcing & $\sim$204\,Hz & Battery (frozen) & Li-ion 6S1P, 108\,Wh, 450\,g \\
Fuselage envelope & $\varnothing$107 $\times$ 533\,mm & Structure & segmented-FDM semi-monocoque \\
\bottomrule
\end{tabular}
\end{table}

\begin{figure}[htbp]
\centering
\begin{subfigure}[t]{0.48\textwidth}
  \centering
  \includegraphics[width=\linewidth]{vbat_render}
  \caption{Rendered vehicle, tail-down on its landing legs.}
\end{subfigure}\hfill
\begin{subfigure}[t]{0.48\textwidth}
  \centering
  \includegraphics[width=\linewidth]{fuselage_layout}
  \caption{Fuselage packaging and bay stack from the layout stage.}
\end{subfigure}
\caption{Pipeline-generated design artifacts: every figure is
regenerated from configuration by the deterministic pipeline; none is
hand-drawn.}
\label{fig:vehicle}
\end{figure}

The commercial-component freeze (pipeline stage 12) selected, against
requirements \emph{derived} from the converged design point, a
Pixhawk-class PX4 flight controller, a telemetry-capable 80\,A-class
ESC, the pinned 3-blade $8\times6$ propeller, sub-10\,g-class
control-surface servos, and a 21700 Li-ion 6S pack that undercuts the
sized LiPo pack by $\sim$203\,g --- flipping the as-selected all-up
mass 182\,g \emph{below} the conceptual closure. The as-selected
re-solve (stages 13--15) then re-solved the affected subsystems with
real hardware masses and envelopes and confirmed physical fit in the
unchanged $\varnothing$107\,mm hull.

Crucially, the pipeline also reports what does \emph{not} close.
Standing findings pinned at the final snapshot include: a marginal ESC
cold plate (its $\sim$37\,W hover load needs a plate heavier than the
ESC mass allocation, with only a few $^{\circ}$C margin); a battery
pack that, at the 40\,$^{\circ}$C design ambient and nominal
90\,m$\Omega$ resistance, ends the mission $\sim$5\,$^{\circ}$C over
its 60\,$^{\circ}$C limit (making measured pack resistance a
procurement acceptance criterion); and avionics-bay and motor mass
allocations exceeded by $\sim$25\,g and $\sim$79\,g respectively.
These are carried openly as accepted risks with named retirement
paths, not tuned away --- the behavior the standing-findings
discipline of Section~\ref{sec:verification} was designed to produce.

\subsection{Design-point evolution}
\label{sec:evolution}

Table~\ref{tab:evolution} traces the design point across the tagged
snapshots. The trajectory is deliberately non-monotonic: the design
got \emph{heavier} whenever a model was made more honest (the
3D-printability penalty at v0.2.0; the battery internal-resistance
accounting at v0.5.1, $+215$\,g) and lighter whenever architecture or
procurement evidence justified it (the externally reviewed
semi-monocoque re-baseline at v0.4.0; the larger, lower-disk-loading
propeller the component freeze forced at v0.4.1). Each row is a
complete, internally consistent, released design snapshot with its
regression pins, CFD references, CAD exports, and changelog updated in
the same change --- there is no intermediate state in the history in
which the artifacts disagree.

\begin{table}[htbp]
\centering
\caption{Design-point evolution across tagged snapshots. Every row is
a fully re-converged, regression-pinned release.}
\label{tab:evolution}
\small
\begin{tabularx}{\textwidth}{@{}llllX@{}}
\toprule
\textbf{Tag} & \textbf{Date} & \textbf{MTOW} & \textbf{Hover} &
\textbf{Driver of the change} \\
\midrule
v0.1.0 & Jul 09 & 2.50\,kg & $\sim$770\,W &
Initial closure: NACA 2412 wing, 195\,mm COTS fan assumption \\
v0.2.0 & Jul 10 & 3.06\,kg & $\sim$1030\,W &
Honest 3D-printed-structure penalty enters the closure; vibration
isolation and modularity split lines added \\
v0.3.0 & Jul 10 & 3.06\,kg & --- &
Thermal paths sized as structure; marginal ESC cold plate first
reported as a standing finding \\
v0.4.0 & Jul 11 & 2.376\,kg & 710\,W &
Semi-monocoque clamshell fuselage from independent expert review;
structural budget re-baselined to an explicit member model \\
v0.4.1 & Jul 15 & 2.302\,kg & 652\,W &
Component freeze reveals the 195\,mm fan class is unbuyable at this
scale; propulsion decision amended by ADR to a 203\,mm 3-blade
propeller in the airframe duct \\
v0.4.5 & Jul 17 & 2.303\,kg & --- &
Li-ion 21700 pack wins the battery freeze; battery mission-transient
thermal check added from external review \\
v0.5.0 & Jul 17 & 2.303\,kg & --- &
Stall-speed inconsistency found by expert review (placeholder
$C_{L,\max}$ survived in one report); wing re-sized to the real
airfoil limit \\
v0.5.1 & Jul 19 & 2.518\,kg & 743\,W &
NACA 4412 re-selection exploits the now-honest stall constraint;
pack $I^2R$ losses folded into closure ($+215$\,g accepted) \\
v0.5.2 & Jul 21 & 2.518\,kg & 743\,W &
Versioned aerodynamic hand-off contract iterated with the external
solver consumer (six schema releases) \\
\bottomrule
\end{tabularx}
\end{table}

The vehicle that emerged is materially different from the one the first
snapshot described, and MTOW alone understates how much moved.
Table~\ref{tab:derived} therefore traces the \emph{derived} quantities
--- the class-(c) and class-(d) parameters of
Section~\ref{sec:taxonomy} --- through the same snapshots, which is
where the coupling of Fig.~\ref{fig:coupling} becomes visible as
program history rather than as a diagram. Four propagation chains
account for essentially the whole trajectory:

\begin{itemize}[nosep]
\item \textbf{Structural honesty drives everything through the
  closure.} Admitting the 3D-printing penalty at v0.2.0 raised the
  structural fraction, and because the fractions multiply MTOW in
  Eq.~(\ref{eq:closure}), the effect compounded: $+0.56$\,kg of MTOW,
  a disk loading rising past 1000\,N/m$^2$
  (Eq.~(\ref{eq:dl}) at the then-fixed 195\,mm rotor), hover power up
  to $\sim$1.03\,kW, hover thrust demand at 39\,N, and a wing grown to
  0.244\,m$^2$ / 1.21\,m span purely because the stall-limited wing
  loading was unchanged while weight was not. Replacing that
  area-density estimate with the explicit member model of the
  externally reviewed semi-monocoque architecture at v0.4.0 ran the
  same chain backwards and recovered more than the whole penalty.
\item \textbf{Rotor diameter moved four disciplines at once.} The
  195\,$\rightarrow$\,203\,mm change at v0.4.1 --- forced by
  procurement, not by analysis --- gave $+8\%$ disk area, dropped disk
  loading 780\,$\rightarrow$\,697\,N/m$^2$, cut hover power
  710\,$\rightarrow$\,652\,W, shrank the pack from 3.71 to 3.50\,Ah,
  and simultaneously retuned the 1/rev vibration forcing
  211\,$\rightarrow$\,204\,Hz, which the isolation stage must be sized
  against. It also moved the guard: the rotor thrust ceiling
  $T_{\max}$ was re-derived for the prop-in-duct class
  (50\,$\rightarrow$\,35\,N), leaving hover at $\sim$84\% of capability
  --- a tighter margin than before the change, on a lighter aircraft.
\item \textbf{$C_{L,\max}$ is the lever on wing geometry, and it moved
  twice.} Until v0.5.0 the stall boundary of Eq.~(\ref{eq:stall}) was
  evaluated with a pre-selection placeholder
  $C_{L,\max}=1.4$ while every other wing number came from the selected
  NACA 2412 --- so the wing was sized to a lift coefficient the airfoil
  could not deliver, and the 12\,m/s stall requirement was silently
  unmet at a true $\sim$12.8\,m/s. Correcting it to the real 1.2186
  \emph{penalised} the wing (loading 123.2\,$\rightarrow$\,107.2\,N/m$^2$,
  area up to 0.211\,m$^2$) without touching MTOW or power. Only then
  did section lift become a genuine lever, and the NACA 4412
  re-selection at v0.5.1 raised $C_{L,\max,3D}$ to 1.4886, taking the
  wing loading to 131.2\,N/m$^2$ --- a wing smaller than the one before
  the error was found, but now honestly sized.
\item \textbf{Pack resistance closed the last inconsistency.} The
  v0.5.1 amendment replaced a constant $\eta_{\mathrm{bat}}=0.97$
  (which implied an unrealistic $\sim$23\,m$\Omega$ pack) with the
  mission-averaged $I^{2}R$ efficiency of the nominal 90\,m$\Omega$
  pack, 0.916, iterated to the closure fixed point --- because the
  heavier battery raises hover current, which lowers the efficiency
  again. The same $R_{\mathrm{pack}}$ now feeds both
  Eq.~(\ref{eq:closure}) and the thermal loads of
  Eqs.~(\ref{eq:thermalss})--(\ref{eq:transient}), which is the
  coupling flagged in Section~\ref{sec:coupling}. The honest accounting
  cost $+215$\,g of MTOW and pushed the wing back out to 0.188\,m$^2$
  at the new weight.
\end{itemize}

One conservatism is worth naming explicitly, because it is visible in
Table~\ref{tab:derived} as an apparent inconsistency. The mass closure
computes cruise power from a configured $L/D=8.0$
(Eq.~(\ref{eq:pcr})), while the wing stage reports the aerodynamic
$L/D$ of the selected polar at the cruise lift coefficient --- 13.35 at
the final snapshot. The latter is deliberately \emph{not} fed back into
the closure: doing so would create a back-edge in the one-way pipeline
of Section~\ref{sec:pipeline}, and the gap is retained as a named
conservatism on cruise energy. Consequently MTOW and hover power do not
move when the wing is re-selected, which is why the v0.5.0 and v0.5.1
wing changes appear in Table~\ref{tab:derived} without a corresponding
mass change.

\begin{table}[htbp]
\centering
\caption{Derived design parameters across the tagged snapshots ---
the quantities that actually characterise the vehicle, none of which is
configured. MTOW, hover power, discharge rate, rotor diameter, forcing
frequency, thrust guard, $C_{L,\max}$, and wing loading are as recorded
in each snapshot's release notes; wing area, span, and MAC follow
Eqs.~(\ref{eq:stall}) and (\ref{eq:winggeom}), disk loading follows
Eq.~(\ref{eq:dl}), and the hover thrust requirement is the weight
times the configured 1.3 thrust margin. A dash marks a quantity not
reported at that snapshot; $^{\dagger}$ marks the pre-selection
placeholder $C_{L,\max}$ corrected at v0.5.0.}
\label{tab:derived}
\footnotesize
\setlength{\tabcolsep}{3.5pt}
\begin{tabular}{@{}lcccccc@{}}
\toprule
& \textbf{v0.1.0} & \textbf{v0.2.0} & \textbf{v0.4.0}
& \textbf{v0.4.1/.4.5} & \textbf{v0.5.0} & \textbf{v0.5.1/.5.2} \\
& Jul 09 & Jul 10 & Jul 11 & Jul 15--17 & Jul 17 & Jul 19--21 \\
\midrule
MTOW [kg]                          & 2.50 & 3.06 & 2.376 & 2.302 & 2.303 & 2.518 \\
Hover power [W]                    & 770  & 1030 & 710   & 652   & 652   & 743 \\
Hover thrust req.\ [N]             & 31.9 & 39.0 & 30.3  & 29.3  & 29.3  & 32.1 \\
Peak discharge [C]                 & 9.0  & ---  & 8.6   & ---   & ---   & 8.1 \\
\addlinespace
Rotor diameter [mm]                & 195  & 195  & 195   & 203   & 203   & 203 \\
Disk loading [N/m$^2$]             & 821  & 1005 & 780   & 697   & 697   & 763 \\
1/rev forcing [Hz]                 & ---  & 211  & 211   & 204   & 204   & 204 \\
$T_{\max}$ guard [N]               & ---  & 50   & 50    & 35    & 35    & 35 \\
\addlinespace
$C_{L,\max,3D}$                    & 1.400$^{\dagger}$ & 1.400$^{\dagger}$
                                   & 1.400$^{\dagger}$ & 1.400$^{\dagger}$
                                   & 1.2186 & 1.4886 \\
Wing loading [N/m$^2$]             & 123.2 & 123.2 & 123.2 & 123.2 & 107.2 & 131.2 \\
Wing area [m$^2$]                  & 0.199 & 0.244 & 0.189 & 0.183 & 0.211 & 0.188 \\
Span [m]                           & 1.09  & 1.21  & 1.06  & 1.05  & 1.12  & 1.06 \\
MAC [m]                            & 0.182 & 0.201 & 0.177 & 0.175 & 0.187 & 0.177 \\
Cruise $L/D$ (polar)               & ---   & ---   & ---   & 13.22 & 12.76 & 13.35 \\
\addlinespace
Airfoil                            & 2412 & 2412 & 2412 & 2412 & 2412 & 4412 \\
\bottomrule
\end{tabular}
\end{table}

\subsection{Process evidence: quantifying the AI's role}
\label{sec:process-results}

Because every AI-implemented commit carries a machine-readable
co-authorship trailer, the division of labor is recoverable directly
from the git history, providing unusually direct process evidence.

\paragraph{Two phases, one repository.} The project history contains a
natural comparison. In a first, entirely manual phase (4~February --
27~March), 28 commits over seven calendar weeks carried the study
through mission sizing, initial take-off-weight iteration, airfoil
selection, and control-vane design --- roughly the first three of the
eventual sixteen pipeline stages, with physics gradually migrating out
of notebooks. In the AI-assisted phase (4--24~July), 109 commits over
three calendar weeks of part-time work delivered the remaining
thirteen pipeline stages plus the full governance and verification
infrastructure: parametric CAD, CFD cases, vibration, thermal, mass
properties, wiring, component freeze, the as-selected re-solve,
sixteen executed notebooks in CI, 147 pinned regression tests, and
thirteen tagged design releases (v0.1.0 through v0.5.2 in twelve
days). While the two phases are not perfectly controlled (the second
benefited from the first's groundwork), the contrast in scope per
calendar week is roughly an order of magnitude and is consistent with
the acceleration claims in the surrogate-model literature, achieved
here at the workflow level rather than inside a single discipline.

\paragraph{Contribution split.} Of the 109 AI-phase commits, 70
(64\%) carry AI co-authorship trailers; the remainder are human-only
commits (merges, reviews, manual fixes, procurement research). All 23
engineering pull requests of the AI phase were human-reviewed before
merge. The final repository comprises $\sim$11{,}200 lines of
deterministic physics library code, $\sim$1{,}800 lines of tests, 18
commented configuration files, 16 orchestration notebooks, and 18
architecture decision records --- the ADRs and changelog entries
drafted by the agent as part of the changes they document.

\paragraph{Error-catching channels.} The history also documents
\emph{which} safety net caught each class of error, summarized in
Table~\ref{tab:errors}. The pattern is consistent with the layered
design of Section~\ref{sec:method}: automated pins caught silent
drift, schema-governed hand-offs caught cross-tool integration faults,
and independent human expertise caught conceptual errors invisible to
any automated check --- confirming that expert review is a
load-bearing element of the methodology, not a formality.

\begin{table}[htbp]
\centering
\caption{Representative faults during the case study and the
methodology layer that caught each.}
\label{tab:errors}
\small
\begin{tabularx}{\textwidth}{@{}XX@{}}
\toprule
\textbf{Fault} & \textbf{Caught by} \\
\midrule
Propulsion/vane CFD cases silently meshing an obsolete design point &
Cross-consumer consistency discipline; cases re-sourced from hand-offs
with a design-point stamp that refuses stale reuse \\
\addlinespace
Stall requirement silently unmet (placeholder $C_{L,\max}$ vs.\
selected airfoil) & Independent expert review (v0.5.0) \\
\addlinespace
Unrealistically optimistic battery efficiency (implied
$\sim$23\,m$\Omega$ pack) & Expert review prompting the transient
model; closure re-derived from pack physics (v0.4.5--v0.5.1) \\
\addlinespace
Wing control points inside the fuselage in the external solver &
Versioned hand-off schema iteration with the consumer (placement
anchor made mandatory from schema 1.5.0) \\
\addlinespace
Unbuyable propulsion assumption (195\,mm fan class) & Deterministic
component-freeze stage; assumption amended by ADR (v0.4.1) \\
\addlinespace
Refactoring risk while moving physics out of notebooks &
Byte-identical hand-off verification (ADR-0013, ADR-0018) \\
\bottomrule
\end{tabularx}
\end{table}

\subsection{Discussion}
\label{sec:discussion}

\paragraph{Answering the research question.} The case study supports
an affirmative, but conditional, answer to the RQ. The AI accelerated
the program by roughly an order of magnitude in delivered scope per
calendar week --- but the acceleration was \emph{not} obtained by
letting the model produce engineering answers. Every number in
Table~\ref{tab:designpoint} traces to deterministic, configured,
regression-tested physics; the AI's leverage came from (i)
implementing multidisciplinary changes completely, so that a design
move lands with its pins, consumers, decision record, and changelog in
one reviewable diff; (ii) eliminating the bookkeeping latency that
normally separates a design decision from a fully consistent design
state; and (iii) making honest-model refinements cheap enough that
unfavorable results (a $+215$\,g mass growth, a marginal cold plate)
were adopted rather than resisted.

\paragraph{Division of labor.} The observed equilibrium assigns the
AI what LLM agents are reliably good at --- exhaustive propagation,
code implementation, documentation, consistency --- and reserves for
humans what they were uniquely necessary for in this program:
requirement setting, acceptance of every decision, procurement
judgment, and conceptual review. It is notable that both of the case
study's most consequential errors (the stall inconsistency and the
optimistic battery efficiency) were latent \emph{before} the AI phase
and were surfaced during it, one by human expert review and one by an
expert-prompted model refinement the AI then implemented. Neither the
AI alone nor the automation alone would have sufficed; the methodology
is precisely the structure that lets each catch what the other misses.

\paragraph{The methodology is a tool, and its outputs are development
cycle time and quality of work.} Two effects should be separated. The
first is \emph{cycle time}. In the case study the interval between a
design directive and a fully consistent, released design state fell
from the order of weeks to the order of a day: thirteen tagged
snapshots in twelve days, nine of which moved the design point
(Table~\ref{tab:evolution}), each internally complete down to its
regression pins, CAD exports, CFD reference values, decision record,
and changelog. That reduction is not
principally a matter of computing faster --- the physics in
Section~\ref{sec:equations} runs in seconds either way --- but of
eliminating the manual bookkeeping that normally separates a decision
from a consistent design state, and which in conventional practice is
what makes engineers reluctant to revisit a converged point at all. The
second effect is \emph{quality of work}, and it is the more important
of the two. Because a change is only accepted when every coupled
quantity in Fig.~\ref{fig:coupling} has been re-derived and every check
in Table~\ref{tab:constraints} re-evaluated, the class of error that
dominates fast-moving programs --- a parameter updated in one place and
not in the three places that depend on it --- is largely removed from
the space of possible outcomes. The case study's own history is the
evidence: the two most consequential errors were both
\emph{inconsistencies} of exactly this kind, and both had survived the
manual phase undetected.

\paragraph{Who benefits, and why the effect is largest for
integration-driven teams.} The framework supplies, as tooling, what
established design houses possess as institutional heritage:
enforced conventions, propagation discipline, and a record of why each
number is what it is. For an experienced team this mainly buys
iteration speed and protection against fatigue-driven slips during a
compressed schedule. For the far larger population of drone developers
whose work is integration-led rather than analysis-led --- assembling
COTS propulsion, avionics, and energy storage into an airframe, often
without a heritage sizing database or a dedicated loads group --- the
benefit is more fundamental. Such teams are exposed precisely where
Fig.~\ref{fig:coupling} is dense and non-obvious: that rotor diameter
is a mass parameter, that pack resistance is a mass parameter, that
$C_{L,\max}$ is a geometry parameter. The component-freeze and
as-selected re-solve stages address the characteristic failure of this
mode of working directly, by re-deriving the design against datasheet
reality instead of against the assumptions the sizing was run with.

\paragraph{Use as an instrument for formal design review.} A further
practical use emerged during the case study, and it may be the most
readily transferable. Because the final pipeline stage is a pure reader
that regenerates the design point, the margins table, the frozen
hardware list, and every standing finding from the hand-offs on every
run, it constitutes a design review package that cannot be stale by
construction. Three properties make it usable as a review instrument
rather than merely a report. Unmet requirements are \emph{structurally
un-hideable}: a negative margin is re-emitted on every run until it is
retired, so the review agenda is generated rather than assembled, and
the incentive to quietly re-tune an inconvenient assumption is removed.
Every number is traceable in one step to the configured parameter and
equation that produced it, so a reviewer's challenge can be answered ---
or the design re-converged under the reviewer's alternative assumption
--- within the review meeting rather than as an action item. And each
decision under review has a dated record stating its context and
consequences, with superseded decisions amended rather than rewritten,
so a reviewer can audit the reasoning trail and not only the current
state. In this program the mechanism was exercised in earnest: external
review findings entered as directives, and the resulting design moves
are traceable from finding to decision record to re-converged snapshot
(Table~\ref{tab:evolution}).

\paragraph{Generality.} Nothing in
Sections~\ref{sec:pipeline}--\ref{sec:iteration} is specific to
tail-sitters: the constitution, one-way pipeline, pins, ADR
discipline, and interaction protocol transfer directly to any
configuration-driven engineering program with computable analysis
stages. The tail-sitter case exercised the framework hard ---
hover/cruise coupling, dual control effectors, thermal and vibration
constraints, and a procurement reality check all interact --- which is
exactly why it was chosen as the validation vehicle.

\paragraph{Limitations.} The study is a single case with a team
already expert in the domain and the tooling; the phase comparison is
suggestive, not controlled. All analyses are conceptual-level
(momentum theory, lifting-line-class aerodynamics, first-order thermal
and structural models), and no hardware has flown --- the PX4-based
flight validation and the external VLM/BEMT and CFD loops are the
designed next steps, wired into the same hand-off contracts.

The most consequential gap is one that the integrative model of
Section~\ref{sec:model} makes visible rather than hides, and it is
marked as such in Fig.~\ref{fig:coupling}: transition is treated here as
an \emph{energy} item --- 40\,s billed at hover power --- but not as a
\emph{sizing constraint}. The tail-sitter literature is consistent that
this is not a safe simplification. Feasible transition depends on
actuator dynamics, rate limits, and time delay as well as on the
conceptual parameters \citep{zhong2021,li2020,banazadeh2016}; the transition corridor
widens with available propulsive power and narrows with weight
\citep{tang2018}; and imposing a minimum-transition-time constraint
inside the sizing loop has been shown to change the converged takeoff
weight by nearly 5\% \citep{cong2024}, with trajectory-coupled
optimisation reporting far larger energy differences
\citep{kaneko2024}. The present closure therefore cannot certify that
the converged vehicle can complete a transition at all, only that it has
the energy to attempt one; at 84\% of the rotor thrust guard, corridor
width is exactly the quantity most likely to bind. Adding a
transition-feasibility stage --- which the framework accommodates
without structural change, as a further pipeline stage with its own
pinned margin --- is the highest-value next increment to the model, and
is a more informative admission than a generic appeal to fidelity: the
coupling map states which box is empty. The ingredients for that
stage exist in the literature: robust transition and re-transition
procedures with explicit corridors \citep{marvakov2021}, a
demonstrated transition method for ducted-fan tail-sitters
\citep{zhang2013}, nondimensional ducted-fan aerodynamic models
suitable for the corridor computation \citep{ohanian2012}, and
guidance on the aerodynamic-model fidelity that computation needs
\citep{anderson2021,arias2025}. Finally, the
methodology's guardrails bound, but cannot eliminate, LLM
fallibility: the protection is structural (nothing merges without
passing pins, CI, and human review), and the residual risk
concentrates exactly where automated checks are weakest ---
conceptual assumptions --- which is why independent expert review
remains mandatory in the loop.
